% ---------------------------------------------------------------- CAPA
\begin{titlepage}
\thispagestyle{empty}
\pdfbookmark[1]{Capa}{capa}
\begin{tikzpicture}[remember picture,overlay]
  % fundo
  \fill[tinta] (current page.south west) rectangle (current page.north east);
  % selo no alto, à direita
  \node[anchor=north east,fill=corC,text=white,rounded corners=3.5mm,inner xsep=4mm,inner ysep=2.2mm,font=\small\bfseries]
    at ($(current page.north east)+(-1.7,-1.9)$) {ENEM 2021–2025 · 79 questões analisadas};
  % título
  \node[anchor=north west,text width=17cm,align=left] at ($(current page.north west)+(1.9,-3.3)$) {%
    {\color{corC}\bfseries\large REVISÃO RÁPIDA · ENSINO MÉDIO}\\[6mm]
    {\color{white}\fontsize{62}{64}\selectfont\bfseries Física\\[1.5mm] no\hspace{0.22em}ENEM}\\[7mm]
    {\color{white}\fontsize{20}{26}\selectfont Os 3 assuntos que mais caem,\\ explicados com exemplos do dia a dia}};
  % ---- ilustração central: bola lançada (mecânica), lâmpada no circuito (eletricidade), onda (ondas)
  \begin{scope}[shift={($(current page.south west)+(1.9,10.2)$)}]
    % chão
    \draw[white!35!tinta,line width=1.2pt] (0,0) -- (17.2,0);
    % trajetória do lançamento
    \draw[corA,line width=2.6pt,dash pattern=on 7pt off 5pt,domain=0:8,samples=60,smooth] plot (\x,{2.6*\x-0.325*\x*\x});
    \foreach \x in {2,4,6}{\fill[corA!60] (\x,{2.6*\x-0.325*\x*\x}) circle (0.13);}
    \fill[white] (0,0) circle (0.2);
    \fill[corA] (8,0) circle (0.26);
    % circuito com lâmpada acesa
    \begin{scope}[shift={(10.6,0.9)}]
      \draw[corB,line width=2.4pt,line join=round] (0,0) -- (0,3.4) -- (5.6,3.4) -- (5.6,0) -- cycle;
      % pilha
      \fill[tinta] (-0.5,1.35) rectangle (0.5,2.05);
      \draw[corB,line width=2.4pt] (-0.55,2.0) -- (0.55,2.0);
      \draw[corB,line width=4pt] (-0.28,1.45) -- (0.28,1.45);
      % lâmpada
      \fill[tinta] (5.6,1.7) circle (0.85);
      \foreach \a in {0,45,...,315}{\draw[ouro!80!white,line width=1.6pt,line cap=round] ($(5.6,1.7)+(\a:0.95)$) -- ($(5.6,1.7)+(\a:1.3)$);}
      \fill[ouro!85!white] (5.6,1.7) circle (0.62);
      \draw[tinta,line width=1.4pt] ($(5.6,1.7)+(-0.3,-0.3)$) -- ($(5.6,1.7)+(0.3,0.3)$);
      \draw[tinta,line width=1.4pt] ($(5.6,1.7)+(-0.3,0.3)$) -- ($(5.6,1.7)+(0.3,-0.3)$);
    \end{scope}
  \end{scope}
  % onda atravessando a página
  \begin{scope}[shift={($(current page.south west)+(0,7.6)$)}]
    \draw[corC,line width=3pt,samples=200,domain=0:21,smooth] plot (\x,{1.05*sin(\x*100)});
    \draw[corC!45,line width=1.5pt,samples=200,domain=0:21,smooth] plot (\x,{0.55*sin(\x*100+60)-0.2});
  \end{scope}
  % três selos dos assuntos
  \foreach \i/\cor/\ic/\nome/\sub in {0/corA/\faCarSide/Mecânica/{movimento, forças e energia}, 1/corB/\faBolt/Eletricidade/{circuitos e conta de luz}, 2/corC/\faWaveSquare/Ondas/{som, luz e fenômenos}}{
    \begin{scope}[shift={($(current page.south west)+(1.9+5.9*\i,3.9)$)}]
      \fill[\cor] (0.8,0.8) circle (0.8);
      \node[white,font=\LARGE] at (0.8,0.8) {\ic};
      \node[anchor=north west,white,font=\bfseries\Large,inner sep=0pt] at (0,-0.35) {\nome};
      \node[anchor=north west,white!80!tinta,font=\normalsize,inner sep=0pt] at (0,-1.05) {\sub};
    \end{scope}}
  % rodapé da capa
  \node[anchor=south west,text width=13.5cm,align=left,white!85!tinta,font=\small] at ($(current page.south west)+(1.9,1.0)$)
    {Resumos curtos, questões reais comentadas, 6 questões de treino com gabarito e a análise completa das provas.};
  \node[anchor=south east,white,font=\bfseries\large] at ($(current page.south east)+(-1.7,1.0)$) {Espaço Cria};
\end{tikzpicture}
\end{titlepage}
\setcounter{page}{2}   % a capa é a página 1: a numeração impressa coincide com a do leitor de PDF
